\section*{5 Codes Closer to the GV Bound}
\subsection*{5.1 The construction}
We use our parity sampler to amplify the distance of a given base code via direct sum lifting. More formally, given$k\in\mathbb N$ and$\varepsilon>0$, let$\varepsilon_0=\varepsilon_0(\varepsilon)$ soon to be determined, and let
\begin{itemize}
\item $C_0:\mathbb F_2^k\to\mathbb F_2^n$ be an$\varepsilon_0$-balanced code, and,
\item $\mathbf W\sim[n]^t$ be an$(\varepsilon_0,\varepsilon)$-parity sampler which is$M$-discretizable.
\end{itemize}
Denote$\bar n=|\operatorname{Supp}(\mathbf W)|\leq M$. Thus, the lifted code$C=\dsum_{\mathbf W}(C_0)\subseteq\mathbb F_2^{\bar n}$ is clearly$\varepsilon$-balanced. For the base code, we use the codes by Guruswami and Indyk, that admit$O_{\varepsilon_0}(n)$-time encoding and decoding.
\paragraph{Theorem30 ([15]).} For every integer$n$ and any$\varepsilon_0>0$ there exists an$\varepsilon_0$-balanced code$C_0\subseteq\mathbb F_2^n$ of rate$\Omega(\varepsilon_0^3)$. Furthermore,$C_0$ is encodable in time$\exp(\operatorname{poly}(1/\varepsilon_0))n$ and decodable from$1/4-\varepsilon_0$ fraction of errors in time$\exp(\operatorname{poly}(1/\varepsilon_0))n$.\footnote{One can get a better randomized encoding time of$\operatorname{poly}(1/\varepsilon_0)$.}
\paragraph{Setting parameters.} Given$\varepsilon>0$, we henceforth set:
\begin{enumerate}
\item The initial bias of$C_0$ to$\varepsilon_0=\frac12\cdot2^{-\sqrt{\log(1/\varepsilon)}}$.
\item The number of steps over$H$ to$t=\lceil\sqrt{\log(1/\varepsilon)}\rceil$.
\item The degree$d$ of$H$ to be the smallest even integer for which
\[c_{\mathrm{spec}}\cdot c_{\mathrm{rand}}\cdot\frac1{\sqrt d}\log\left(\frac{\sqrt d}{c_{\mathrm{rand}}}\right)\leq\varepsilon_0.\]
This is chosen so that Corollary21 gives, with high probability, a symmetric hypergraph$H$ such that, by Corollary26,$\mathbf W_{H,t}$ is a$(\varepsilon_0,(2\varepsilon_0)^t)$-parity sampler.
\end{enumerate}
Using the above parameters in Theorem25, together with the random hypergraph of Corollary21. we get the following (randomized) error correcting code$C$ [Theorem31, stated above].
\begin{proof}[Proof of Theorem31]
Given$k\in\mathbb N$ and$\varepsilon>0$, we set$\varepsilon_0,t$, and$d$ as above. Theorem30 gives us a code$C_0:\mathbb F_2^k\to\mathbb F_2^n$ which is$\varepsilon_0$-balanced, for$n=O(k/\varepsilon_0^3)$. By Corollary21 we can output a$\lambda=c_{\mathrm{rand}}/\sqrt d$-mixing$H$ with probability at least$1-2^{-n}$ in polynomial time. This is our preprocessing step. By Corollary26 and our choice of parameters, indeed
\[\left(\varepsilon_0+c_{\mathrm{spec}}\cdot\lambda\log\frac1\lambda\right)^t\leq(2\varepsilon_0)^t\leq\varepsilon,\]
so$\mathbf W=\mathbf W_{H,t}$ appropriately amplifies the distance. Moreover,$\mathbf W$ is$M$-discretizable for
\begin{align*}
M&=nd^t=O\left(\frac k{\varepsilon_0^3}\right)\cdot d^{\lceil\sqrt{\log(1/\varepsilon)}\rceil}\\
&=k\cdot2^{(3+\log d)\sqrt{\log(1/\varepsilon)}+\log d+O(1)}.
\end{align*}
We proceed to bounding$\log d$ in terms of$\varepsilon$. Recalling how we set$d$, we see that$\frac c{\sqrt d}\log d=2^{-\sqrt{\log(1/\varepsilon)}}$, where$c$ is some universal positive constant. From this we can infer that
\[d\leq c^2 2^{2\sqrt{\log(1/\varepsilon)}}\cdot\log^3\left(c\cdot2^{\sqrt{\log(1/\varepsilon)}}\right),\tag{7}\]
so$\log d=2\sqrt{\log(1/\varepsilon)}+O(\log\log(1/\varepsilon))$. (We are assuming that$\varepsilon$ is smaller than some small constant, which we can without loss of generality.) Hence,
\[M=\frac k{\varepsilon^2}\cdot2^{O(\sqrt{\log(1/\varepsilon)}\cdot\log\log(1/\varepsilon))}.\]
Since$\mathbf W$ is$M$-discretizable,$C=\dsum_{\mathbf W}(C_0)$ has block length$\bar n\leq M$, as required. Drawing$H$ at random, by Corollary21 takes
\[O(dn)=\widetilde O\left(\frac1{\varepsilon_0^3}\right)\cdot k=\exp(\sqrt{\log(1/\varepsilon)})\cdot k\]
time. Given the hypergraph$H$, the encoding amounts to computing$C_0(x)$ and taking parities of its coordinates according to$\mathbf W$. This takes$\exp(\operatorname{poly}(1/\varepsilon_0))\cdot k+M\cdot t$ time, following Theorem30.
For the ``moreover'' part, one can verify that we can tolerate a$\operatorname{polylog}(d)$ multiplicative factor in$\lambda$ with no substantial loss in parameters.
\end{proof}
